\subsection{滑らかな先頭三次を持つ非対称alpha包絡の全中央三次}
\label{subsec:cubic-CY-three-alpha-smooth-boundary}

\begin{theorem}\label{thm:cubic-CY-three-alpha-smooth-boundary}
三交換子の $\alpha$ 型で $\gamma\ne0$、先頭三次曲線が滑らかとする。
中央一次元を $kt$、$S=\mathcal B/(t)$ とする。
Hesse座標で $S$ の関係を
\[
 p yz+qzy+r x^2=0,\qquad\text{巡回}
\]
と書き、$E$ を
\[
 (p^3+q^3+r^3)xyz-pqr(x^3+y^3+z^3)=0
\]
で定める。すると $E$ は滑らかな楕円曲線または三角形である。
どちらの場合も $Z(S)_1=Z(S)_2=0$、$Z(S)_3=kg$ であり、
$g$ を持ち上げる $L\in Z(\mathcal B)_3$ が存在する。
従って全中央三次は $kL+kt^3$ である。

$E$ が滑らかな場合、全ての $L+\lambda t^3$ は整域と内在的rank-oneを与え、
H3条件はアフィン環の有限大域次元性に等しい。
境界平行移動が有限位数なら、中心と局所McKay条件による既存の必要十分判定が適用できる。
$E$ が三角形の場合、先頭包絡 $S$ は
$yz=h x^2,\ zx=h y^2,\ xy=h z^2$ の形へ移り、
$h\ne0$、$h^3\ne1$ となる。
この三角形では境界自己同型が三成分を巡回する。
\end{theorem}
\begin{proof}
滑らかな平面三次曲線のHesse標準形\cite{ADHesse}と
固定出典\cite{AS4CYpinned}の $\alpha$ の必要表示を使う。
$F=s(x^3+y^3+z^3)+vxyz$ とすると $r=-3s$、
規格化 $\gamma=1$ では $p=1-v/2,q=-1-v/2$ である。
特に
\[
 r\ne0,\qquad p-q\ne0,\qquad (p+q)^3+r^3\ne0.
\]
先頭三関係の正則性から、$p,q,r$ の少なくとも二つが非零で、
$p^3=q^3=r^3$ ではない。
逆にこの二条件でIyudu--Shkarin
\cite[定理1.1および系6.1]{ISSklyanin}のquadratic Koszul CY AS3となる。
$t$ は元の整域で正則であり、活動次数の随伴環は
$S[t]$ への表示とHilbert値の比較で同定される。

点行列の行列式が表示した三次式である。
$pqr\ne0$ のとき、その曲線が滑らかでない条件は
$(p^3+q^3+r^3)^3=27p^3q^3r^3$ であり、Hesseの三角形である。
$pq=0$ のときも、$p^3+q^3+r^3\ne0$ は上の滑らかさ条件から従うので座標三角形である。
$pq\ne0$ の三角形を考え、$r$ に三乗根を掛けて
$p^3+q^3+r^3=3pqr$ としてよい。
この差は
$(p+q+r)(p+\zeta q+\zeta^2r)(p+\zeta^2 q+\zeta r)$ であり、
最初の因子は零でない。
Iyudu--Shkarinの三変数Fourier変換は
\[
 (p,q,r)\longmapsto(\zeta^2p+\zeta q+r,\,
                     \zeta p+\zeta^2q+r,\,p+q+r)
\]
を与えるので、一方の交差係数が零となる。
変数の交換と尺度を用いて $p=1,q=0,r=-h$ とできる。
滑らかさの条件は $h^3\ne1$ となる。

滑らかな $E$ では点行列の階数が2であり、
二重線型化は $E$ の自己同型 $\sigma$ のグラフとなる。
固定点なら各関係の可換化が零となるが、これは $F$ の勾配が零の点で、
$F$ の滑らかさに反する。従って $\sigma$ は平行移動である。
変曲点 $o=(1:-1:0)$ を原点に取ると、その像は $(p:q:r)$。
Hesse曲線の3等分点は九つの変曲点で全て一座標が零なので、
$pqr\ne0$ から $\sigma^3\ne1$ である。
逆元は $x,y$ の交換であり、$r\ne0,p\ne q$ から $\sigma^2\ne1$ である。
$\mathcal O_E(1)$ と $\sigma$ による引戻しの差はこの平行移動の3倍に対応し、
同型でない。基点自由pencilの完全列で次数二の積が全射となり、
高次数はテンソル積の正次数で全射となる。
従って $S\to B(E,\mathcal O_E(1),\sigma)$ は全射である。

三角形の表示では、座標直線の開部分で点の後続写像は
\[
 (x,y,0)\longmapsto (y,0,hx)
 \longmapsto (0,h^2 x,y)
 \longmapsto (h^3x,y,0)
\]
となる。各直線の写像は頂点で一致するので三角曲線上の自己同型である。
従って $\sigma,\sigma^2$ は成分を巡回し、
$\sigma^3$ はどの成分上でも恒等にならない。
二重線型関係から $S\to B(E,\mathcal O_E(1),\sigma)$ を得る。
二次の六語の三辺での評価は単元小行列式を持ち階数6。
各成分上の次数1を保つので、次数 $n\ge2$ の積の全射性は
定理\ref{thm:cubic-CY-three-alpha-additive-H3}のpencilの議論で従う。
よってこの場合もTHCRへの全射を得る。

いずれの場合も、このTHCRでは次数1,2,3の中心が零である。
実際、中心斉次切断の非零成分で一次切断二つとの可換性を比に取ると、
$\sigma^m$ がその成分を点ごとに固定することが必要である。
次数 $m=1,2,3$ にそのような成分はない。
全射の核は二次まで零、三次ではHilbert値の差から一次元である。
従って $Z(S)_1=Z(S)_2=0$、$\dim Z(S)_3\le1$。

$t=1$ の包絡は、先頭が今のCY3ポテンシャルで、
残りが次数2以下の巡回ポテンシャルの代数である。
Etingof--Ginzburg\cite[定理3.3.2]{EGdelpezzo}はこれに非スカラーな
濾過次数高々3の中央元を与える。
先頭中心の次数1,2が零なのでその濾過次数は3で、最高記号は非零である。
斉次化して $L$ を得る。複素数体から一般の標数零代数閉体へは、
次数3,4の商と $Z(\mathcal B)_3\to S_3$ の階数を使う。
係数の有限生成体を複素数体に埋め込んでもこれらの階数は変わらず、
像の非零性が移る。$L$ の像 $g$ が $Z(S)_3$ を生成する。
中心の次数1,2の同様の比較と $t$ の正則性により
$Z(\mathcal B)_3=kL+kt^3$ を得る。
$g$ はAS3整域 $S$ で正則なので、Hilbert値比較から
$S/(g)=B(E,\mathcal O_E(1),\sigma)$ である。

滑らかな $E$ では従って良い楕円中心拡大の仮定がすべて成り立ち、
定理\ref{thm:r13-exact-cut-reduction}と
系\ref{cor:normal-cubic-finite-order-admissibility}を適用する。
滑らかな先頭式の元の包絡同型全体は固定出典の変曲点枠の手順で得られ、
$\phi(t)=a t'$、$\phi(L)=bL'+c(t')^3$ と書けば、
切断値の作用は $(c+a^3\lambda)/b$ である。
\end{proof}

公開検証器で行列式、三角形への係数変換、三辺の後続写像と次数二評価を再現する。
\par\noindent\texttt{verify\_cubic\_CY\_three\_alpha\_smooth\_boundary.py}
\par
滑らかな楕円境界の無限位数でのアフィン判定、
およびその幾何学的成分の全判定は残る。

\begin{proposition}\label{prop:cubic-CY-three-alpha-elliptic-character-parity}
滑らかな楕円境界の切断がH3なら、
そのアフィン環 $U$ の任意の有限次元加群 $M$ に対して
$\dim_k\operatorname{Ext}^1_U(M,M)$ は偶数である。
特に一次元表現の点で、三アフィン関係と切断式を可換化した
四式のJacobian行列が偶数階なら、その切断はH3から除外される。
\end{proposition}
\begin{proof}
アフィン包絡 $D=\mathcal B[t^{-1}]_0$ はCY3であり、
$U=D/(\omega+\lambda)$ は中心正則商である。
実際、$\mathcal B$ はCY4で、中央局所化
$\mathcal B[t^{-1}]=D[t,t^{-1}]$ は双加群滑らかさとCY双対を保つ。
定理\ref{thm:r13-split-all-cuts}のLaurent対角の基底変更で
$D$ は双加群滑らかであり、中心正則式 $t-1$ による双対移行でCY3となる。
有限次元右 $D$ 加群 $M$ に対し、CY双加群解消を $M$ とテンソルして双対化すると
$\mathbf R\operatorname{Hom}_D(M,D)\simeq M^*[-3]$ である。
中心正則商の変換で、有限次元 $U$ 加群には
$\mathbf R\operatorname{Hom}_U(M,U)\simeq M^*[-2]$ を得る。

H3なら $U$ は有限大域次元を持つため、有限次元加群は完全である。
従って上の自然双対と完全加群の評価から、
有限次元コホモロジーの複体に対してSerre双対 $[2]$ を得る。
この双対での積の跡は、コチェインの評価で次数 $i,j$ の射を
循環させるとき $(-1)^{ij}$ を生じる。
よって $\operatorname{Ext}^1_U(M,M)$ 上のYoneda積と跡の対合は
非退化な交代形式である。

一次元表現での自己拡大は、上三角行列の非対角成分、
すなわち双対数への表現で記述できる。
その線型方程式は可換化した四式のJacobianであり、
一次元表現では内的微分は零である。
従って自己 $\operatorname{Ext}^1$ の次元は $3$ からその階数を引いた値となる。
\end{proof}

\begin{theorem}[一次元表現の自己拡大と境界楕円曲線]
\label{thm:cubic-CY-three-alpha-elliptic-character-boundary}
滑らかな楕円境界を持つ切断のアフィン環 $U$ の一次元表現を取る。
その点における三二次関係のJacobianを $A$ とし、
関係の符号を揃えて対称行列とする。
一次元加群の自己 $\operatorname{Ext}^1$ の次元は
\[
\begin{array}{c|c}
\text{条件}&\dim\operatorname{Ext}^1\\ \hline
\operatorname{rank} A=0&3\\
\operatorname{rank} A=1&2\\
\operatorname{rank} A=2,\ \mathbb P(\ker A)\in E&1\\
\operatorname{rank} A=2,\ \mathbb P(\ker A)\notin E&0\\
\operatorname{rank} A=3&0
\end{array}
\]
である。従って階数零、または階数2で核点が境界 $E$ 上にある一次元表現は
H3を排除する。
\end{theorem}
\begin{proof}
一次元表現を原点へ移すと、元の三関係の定数項は零となる。
中央切断の定数項も零である。
先頭のHesse係数は $pqr\ne0$、$p\ne q$ なので $p=1$ と規格化する。
先頭の三二次関係の右辺の線型係数を
\[
 A=\begin{pmatrix}a&b&d\\b&c&e\\d&e&f\end{pmatrix}
\]
とする。中央三次を $G$ として、$t^2$ 係数の線型形式を $\ell$ と書く。
次の正規化を使う。
\[
 G\bmod t
 =r(r^3-1)y^3+(r^3-q^3)yzx-q(r^3-1)zyx+r(q^3-r^3)z^3 .
\]
この元は全許容係数で非零である。実際、重み $(0,1,2)$ の
対角三乗根作用の不変次数三には、三つの純三乗と六つの順列語がある。
二次関係の左・右一次倍はその不変関係の全てを与える。
$r\ne0$ で純三乗を消し、さらに二つの独立な差の式で
$xyz,zxy$ を消すと $xzy,yzx,zyx,yxz$ が基底となる。
上の元の $zyx,yxz$ の係数はともに $-q(q^3-1)$、
$xzy$ の係数は $q(q^3-r^3)$ である。
これらが全て零なら $q^3=r^3=1$ となり正則性に反する。

この規格化で中央性の線型方程式を解くと、
持上げの全係数の分母は $r(q-1)$ の因子に限られる。
例えば階数2で核点を $(u:v:1)$ とすると、
$d=-au-bv,\ e=-bu-cv,\ f=au^2+2buv+cv^2$ となり、
厳密な係数比較により
\[
 \ell(u,v,1)=-\frac{ac-b^2}{q-1}
 \bigl((1+q^3+r^3)uv-qr(u^3+v^3+1)\bigr).
\]
同じ計算で、$A$ が階数1なら $\ell$ は $A$ の行空間に入り、
$A=0$ なら $\ell=0$ を得る。
\path{verify_cubic_CY_three_alpha_elliptic_character_boundary.py}
は八変数の有理関数体で中央性の核を計算し、
規格化した全係数、分母、階数2の収縮式、階数1の外積の消滅を確認する。
商の次数三、四の次元は全許容包絡で一定である。
従って中央性は、表示された分母が非零なら特殊化で保たれる。
切断の線型項はこの持上げの $\ell$ の非零倍であり、
定数切断値はJacobianに寄与しない。

階数2で $\ker A$ の第三座標が非零なら $ac-b^2\ne0$ である。
これは余因子行列が、核を表す列と行の積の非零倍だからである。
第三座標が零なら巡回置換で別の座標を第三座標に取れる。
従って階数2で $\ell$ が核を消す条件は正確に核点が $E$ 上にあることとなる。
他の階数では上の行空間条件を使い、
三関係に切断の一行を加えたJacobianの階数から表の値を得る。
\end{proof}

\begin{example}\label{ex:cubic-CY-three-alpha-elliptic-character-obstruction}
三次の先頭関係を $(p,q,r)=(1,2,3)$ とし、$a,b\ne0$ とする。
次の非分裂包絡では、下の中央三次切断は無限位数の滑らかな楕円境界を持ち、
H3から除外される。
\[
\begin{aligned}
 yz+2zy+3x^2&=at(x+y),\\
 zx+2xz+3y^2&=at(x+y),\\
 xy+2yx+3z^2&=btz .
\end{aligned}
\]
$t$ は中心であり、切断は $G=0$ とする。ただし
\[
\begin{aligned}
 G={}&13y^3+\tfrac{19}{6}yzx-\tfrac{26}{3}zyx-\tfrac{19}{2}z^3\\
 &+t\bigl((3b-\tfrac{25}{6}a)yx-\tfrac12ay^2-\tfrac{25}{9}ayz
        -\tfrac32azx-\tfrac{23}{9}azy+\tfrac{23b-27a}{2}z^2\bigr)\\
 &+t^2\bigl((ab-\tfrac{31}{18}a^2)(x+y)
              +(\tfrac{11}{2}ab-3b^2)z\bigr).
\end{aligned}
\]
係数行列が対称なので、この包絡は
補題\ref{lem:sklyanin-central-PBW-CY}により正則CYである。
線型項が先頭式の平行移動から来るなら、$x^2$ から生じる
$x$ の係数 $a$ に対し、$yz+2zy$ から生じる $y,z$ 方向の
相互係数は $a/2$ でなければならない。ここでは零なので非分裂である。
$G$ の四交換子は次数四の元関係の像で零であり、
その最高記号は次数三の商で非零である。
\path{verify_cubic_CY_three_alpha_elliptic_character.py}
がこれらを二変数有理関数体で確認する。
次数三、四のHilbert値は全包絡で一定なので、階数小行列式により中央性は全係数へ特殊化する。

境界は $x^3+y^3+z^3-6xyz=0$、平行移動の点は $(1:2:3)$ である。
$5,17$ で良い還元を持ち、点数はそれぞれ $9,21$。
両還元での素数外torsionの単射により、有理torsionの位数は3を割る。
この点は三座標が非零で変曲点ではなく、従って無限位数である。

$t=1$ の原点は一次元表現である。
三二次関係のJacobianの行空間は $(1,1,0)$、$(0,0,1)$ で張られ、
$G$ の線型項も同じ空間にある。
従って命題\ref{prop:cubic-CY-three-alpha-elliptic-character-parity}により排除される。
\end{example}

\begin{corollary}\label{cor:cubic-CY-three-alpha-cyclic-triangle-excluded}
定理\ref{thm:cubic-CY-three-alpha-smooth-boundary}のうち
境界が成分巡回三角形になる包絡では、全中央三次切断が非整域である。
\end{corollary}
\begin{proof}
定理の最後の座標を用いる。
低次ポテンシャルの二次項は対称二次形式なので、元の三関係は、巡回順で
\[
 yz=h x^2+t(a_{xx}x+a_{xy}y+a_{xz}z)+b_x t^2,qquad
 a_{ij}=a_{ji}
\]
と書ける。$X=x-a_{yz}t$、$Y=y-a_{xz}t$、$Z=z-a_{xy}t$ と置き、
各関係の左辺の混合一次項を移項すると、
\[
 YZ=f_x(X,t),\qquad ZX=f_y(Y,t),\qquad XY=f_z(Z,t)
\]
となる。例えば
$f_x(X,t)=h(X+a_{yz}t)^2+a_{xx}t(X+a_{yz}t)
 +(b_x+a_{xy}a_{xz})t^2$ である。
結合則を $XYZ,YZX$ に適用すると
$Xf_x(X,t)=Yf_y(Y,t)=Zf_z(Z,t)$。
この元は各生成元と可換で、$t=0$ で $h x^3\ne0$ に移る。
従って全中央三次はこれと $t^3$ で張られる。
$X,t$ は独立な可換元なので、全切断は可換二元三次として因数分解する。
一次因子と残る二次因子は整域の包絡で非零であり、
次数三のイデアルで割ってもそれぞれ非零である。
従って全中央三次商に零因子がある。
\end{proof}
